$\bigl[\textbf{КМО 2026, сеньоры, шестая задача.}\bigr]$
В треугольнике $A B C$ угол $B A C$ равен $30^{\circ}$. Пусть $N$ - центр окружности, проходящей через середины отрезков $A B, B C$ и $C A$. Докажите, что $\angle B N C=90^{\circ}$.

\textbf{Решение.} Пусть $O$ — центр описанной окружности треугольника $ABC$. Поместим эту окружность в комплексную плоскость как единичную, а аффиксы точек $A,B,C$ обозначим через $a,b,c$, так что $|a|=|b|=|c|=1$. Поворотом плоскости можно считать, что $a=1$.

Так как $\angle BAC=30^\circ$, то дуга $BC$, не содержащая точку $A$, имеет величину $60^\circ$, а значит, существует комплексное число $u$ с $|u|=1$ такое, что $b=u e^{i\pi/6}$ и $c=u e^{-i\pi/6}$.

Не безызвестно, что комплексные координаты центра окружности Эйлера при единичной описанной окружности равны
\[
n=\frac{a+b+c}{2}.
\]

Читатель может проверить этот факт самостоятельно.

Теперь вычислим векторы $\overrightarrow{NB}$ и $\overrightarrow{NC}$. Имеем
\[
b-n=b-\frac{1+b+c}{2}=\frac{b-1-c}{2},\qquad
c-n=c-\frac{1+b+c}{2}=\frac{c-1-b}{2}.
\]
Подставляя $b=u e^{i\pi/6}$ и $c=u e^{-i\pi/6}$, получаем
\[
b-n=\frac{u(e^{i\pi/6}-e^{-i\pi/6})-1}{2}=\frac{iu-1}{2},\qquad
c-n=\frac{u(e^{-i\pi/6}-e^{i\pi/6})-1}{2}=\frac{-iu-1}{2},
\]
так как $e^{i\pi/6}-e^{-i\pi/6}=2i\sin\frac{\pi}{6}=i$.

Остаётся доказать, что $\overrightarrow{NB}\perp\overrightarrow{NC}$. В комплексной записи это равносильно тому, что число $(b-n)\overline{(c-n)}$ чисто мнимо. Действительно,
\[
(b-n)\overline{(c-n)}=\frac14(iu-1)\,\overline{(-iu-1)}=\frac14(iu-1)(i\overline{u}-1),
\]
поскольку $\overline{u}=u^{-1}$ и потому $\overline{(-iu-1)}=i\overline{u}-1$. Раскрывая скобки, получаем
\[
(iu-1)(i\overline{u}-1)=i^2u\overline{u}-iu-i\overline{u}+1=-1-iu-i\overline{u}+1=-i(u+\overline{u}).
\]
Но $u+\overline{u}=2\operatorname{Re}u$ — вещественное число, значит, $-i(u+\overline{u})$ чисто мнимо. Следовательно, $\operatorname{Re}\bigl((b-n)\overline{(c-n)}\bigr)=0$, откуда $\overrightarrow{NB}\perp\overrightarrow{NC}$.

Итак, $\angle BNC=90^\circ$, что и требовалось доказать.